\chapter{AI 模型、张量与推理图}

本章从零建立推理需要的 AI 基础。我们不讨论训练算法的全部细节，而是聚焦推理引擎必须知道的概念：张量是什么，模型参数是什么，激活是什么，推理图怎样表达计算，算子接口如何连接数据。只有这些概念清楚，后面的 CPU 算子开发才不会变成“对着公式写循环”。

\section{模型不是一个函数名}

一个 AI 模型可以理解为一组参数和一张计算图。参数来自训练结果，例如权重矩阵、卷积核、归一化权重、Embedding 表。推理时输入数据进入计算图，经过一系列算子，产生输出。模型文件保存参数、结构和元数据；推理引擎负责把它加载到内存，并按正确顺序执行。

LCQI 第一阶段的模型很小：两层 MLP。它仍然具备完整模型的基本形态：第一层权重和 bias 是参数，输入向量是输入激活，隐藏层输出是中间激活，第二层输出是 logits。分类结果取 logits 最大值的下标。后续无论模型变成 CNN 还是 Transformer，本质仍然是参数加计算图。

\section{张量、shape、stride 和 dtype}

\term{张量}{tensor} 是带形状的多维数组。shape 描述每一维长度，dtype 描述元素类型，stride 描述每一维移动一步时地址增加多少元素。连续行主序张量的 stride 可以从 shape 推出；切片、转置和某些布局转换会产生非连续 stride。推理引擎必须明确自己支持哪些布局，否则算子读错数据很难排查。

\begin{lstlisting}[language=C++,caption={教学版张量视图：shape 和连续数据}]
#include <cstdint>
#include <span>
#include <vector>

struct TensorView {
    std::span<const float> data;
    std::vector<std::int32_t> shape;
};
\end{lstlisting}

这个结构只表达只读 float 张量视图。真实引擎还要支持可写张量、不同 dtype、stride、内存所有权和设备位置。本册从简单结构开始，是为了让读者先看懂算子接口：算子接收张量，检查 shape，按约定布局读写数据。

\section{参数和激活}

参数是模型加载后长期存在的数据，推理过程中通常只读。激活是某次推理产生的中间数据，生命周期较短。这个区分影响内存管理和并发。参数可以被多个请求和线程共享；激活通常属于一次请求或一个 batch；临时缓冲区可以复用，但不能跨并发请求错误共享。

量化模型中，参数还包含 scale、zero point、分组大小等元数据。一个 int8 权重数组没有 scale 就没有真实数值意义。推理引擎的模型结构必须把量化数据和量化元数据绑在一起，不能只传裸指针。

\section{推理图}

推理图由节点和边组成。节点是算子，边是张量。线性层、ReLU、Softmax、卷积、归一化、Attention 都是节点。推理执行器按依赖顺序执行节点，管理每个边对应的内存。教学项目可以手写固定顺序；更完整的引擎会解析模型文件，构建图，再做内存规划和算子调度。

\begin{mentalmodel}
推理图是数据流，不是代码流。代码流描述程序怎样写，数据流描述张量怎样从一个算子变成下一个算子的输入。算子开发必须同时尊重这两条线。
\end{mentalmodel}

\section{正确性标准}

推理引擎必须有正确性标准。对于 float 算子，常用绝对误差和相对误差；对于量化算子，要比较反量化结果或最终输出类别；对于随机输入，要固定种子；对于模型文件，要有格式校验。没有正确性测试的性能优化没有意义，因为它可能只是算错了。

\begin{exercisebox}
把 LCQI 的两层 MLP 画成推理图，标出每条边的 shape 和 dtype。说明哪些数据是参数，哪些数据是激活，哪些数据是临时缓冲区。
\end{exercisebox}
